\documentclass{article}
\usepackage{graphicx} % Required for inserting images
\usepackage{amsmath,amssymb}
\usepackage{algorithm}
\usepackage{algpseudocode}

\title{federated typiclust}
\author{Ohad Shimoni}
\date{May 2026}

\begin{document}
\maketitle


\section{Federated TypiClust Pipeline}

We propose a federated active learning pipeline inspired by TypiClust~\cite{hacohen2022active}, but designed so that the raw data and most point-level information remain local to the clients. The pipeline is modular: standard cryptographic or federated-learning components are treated as black boxes, while our main contribution is the axis-wise sorted-value mechanism for computing an exact within-cluster $L_1$ typicality score.

Assume there are $K$ clients. Client $k$ holds a private unlabeled dataset
\[
    \mathcal{D}_k = \{x_{k,1},\dots,x_{k,n_k}\}.
\]
The global dataset is
\[
    \mathcal{D} = \bigcup_{k=1}^K \mathcal{D}_k,
\]
but this dataset is never centralized.

\subsection{Step 1: Shared Embedding into $\mathbb{R}^d$}

The first black box provides a shared embedding function
\[
    f : \mathcal{X} \rightarrow \mathbb{R}^d.
\]
Each client applies this embedding locally:
\[
    z_{k,i} = f(x_{k,i}) \in \mathbb{R}^d.
\]

The important requirement is that all clients use the same embedding space. This can be implemented using a public pretrained model, federated supervised learning, or federated self-supervised learning. For example, FedAvg~\cite{mcmahan2017communication} provides a standard black-box method for learning a shared model from decentralized data, and SimCLR~\cite{chen2020simple} provides a standard contrastive method for learning image embeddings. In this work, we treat the construction of $f$ as a black box.

\subsection{Step 2: Private Clustering}

The second black box performs clustering in the shared embedding space. It assigns each embedded point to a cluster:
\[
    c(z_{k,i}) \in \{1,\dots,m\}.
\]
The clustering procedure is assumed to be privacy-preserving. Its output is that each client knows the cluster assignment of its own points, and the protocol can refer to the global clusters
\[
    C_1,\dots,C_m
\]
without collecting all embeddings in one central location.

This step can be implemented using existing privacy-preserving distributed clustering protocols, such as privacy-preserving distributed $k$-means~\cite{vaidya2003privacy,jagannathan2005privacy}. We therefore treat private clustering as a black box.

\subsection{Step 3: Federated Axis-Wise Sorted Values}

For each client $k$, cluster $c$, and coordinate axis $\ell \in \{1,\dots,d\}$, client $k$ considers the local embedded coordinate values
\[
    V_{k,c,\ell}
    =
    \left\{
    z_{k,i}^{(\ell)}
    :
    c(z_{k,i}) = c
    \right\}.
\]

The client sorts these values locally:
\[
    \mathrm{sort}(V_{k,c,\ell}).
\]

The sorted coordinate values are then aggregated across clients so that, for every cluster $c$ and coordinate $\ell$, the global sorted list is
\[
    V_{c,\ell}
    =
    \mathrm{sort}
    \left(
    \bigcup_{k=1}^K V_{k,c,\ell}
    \right).
\]

Thus, the protocol does not reveal complete embedded vectors. Instead, for each cluster, it releases only the sorted marginal coordinate values
\[
    \{V_{c,1},V_{c,2},\dots,V_{c,d}\}.
\]

The server sends these global sorted coordinate lists back to the clients.

\paragraph{Privacy intuition.}
The server sees the coordinate values inside each cluster, but it does not receive the full vectors
\[
    z_{k,i} = (z_{k,i}^{(1)},\dots,z_{k,i}^{(d)}).
\]
In particular, the server does not know which value on axis $1$ belongs to which value on axis $2$, and so on. Therefore, the released information is weaker than releasing all embeddings, but stronger than releasing only coarse histograms. This is the central privacy-utility tradeoff of our method.

\subsection{Step 4: Exact Local Computation of $L_1$ Typicality}

Each client now computes the typicality of each of its local points using the global sorted coordinate lists.

For a point $z$ assigned to cluster $c$, the centralized TypiClust-style $L_1$ score is
\[
    S(z)
    =
    \sum_{z' \in C_c}
    \|z-z'\|_1.
\]

Using the decomposition of the $L_1$ distance,
\[
    \|z-z'\|_1
    =
    \sum_{\ell=1}^d
    |z^{(\ell)} - z'^{(\ell)}|,
\]
we get
\[
    S(z)
    =
    \sum_{\ell=1}^d
    \sum_{v \in V_{c,\ell}}
    |z^{(\ell)} - v|.
\]

Therefore, once the clients receive the global sorted lists $V_{c,\ell}$, each client can compute the exact total $L_1$ distance from its own point to all points in the same cluster, without receiving the full vectors of the other clients.

A point is more typical if its total distance to the cluster is smaller. Therefore, the most typical point in cluster $c$ is
\[
    z_c^\star
    =
    \arg\min_{z \in C_c} S(z).
\]

\paragraph{Efficient computation.}
Let
\[
    V_{c,\ell} = (v_1,\dots,v_n)
\]
be sorted increasingly. For a coordinate value $a=z^{(\ell)}$, suppose $r$ values are smaller than or equal to $a$. Then
\[
    \sum_{j=1}^n |a-v_j|
    =
    a r - \sum_{j=1}^r v_j
    +
    \sum_{j=r+1}^n v_j
    -
    a(n-r).
\]
Using prefix sums over the sorted list, each coordinate contribution can be computed efficiently after a binary search for $r$.

\subsection{Step 5: Secure Argmin for Representative Selection}

For each cluster $c$, the goal is to select the point with minimum total within-cluster $L_1$ distance:
\[
    z_c^\star
    =
    \arg\min_{z \in C_c} S(z).
\]

Each client computes its best local candidate in every cluster it participates in:
\[
    z_{k,c}^\star
    =
    \arg\min_{z \in C_c \cap \mathcal{D}_k} S(z).
\]

The clients and server then run a secure argmin protocol to select the global winner:
\[
    z_c^\star
    =
    \arg\min_k S(z_{k,c}^\star).
\]

This step can be implemented as a standard secure multi-party computation task using secure comparison circuits. General secure computation protocols such as Yao's garbled circuits~\cite{yao1982protocols} and the GMW protocol~\cite{goldreich1987play} provide black-box foundations for such secure comparison and argmin operations.

The protocol reveals only the selected representative for each cluster, or the identity of the client holding it, but does not reveal the scores of all candidates.

\subsection{Step 6: Labeling and Federated Nearest Neighbor}

The selected representatives are sent to a labeling oracle. After labeling, each selected representative has a label:
\[
    (z_c^\star, y_c^\star).
\]

Unlike centralized TypiClust, we do not assume that all labeled representatives must be centralized. Instead, the labeled representatives may remain distributed across the clients that own them.

For prediction, we use a federated nearest-neighbor classifier. Given a query point $x$, the query holder computes
\[
    z = f(x).
\]
The system then finds the closest labeled representative under the $L_1$ metric:
\[
    c^\star
    =
    \arg\min_c
    \|z-z_c^\star\|_1.
\]
The predicted label is
\[
    \hat{y}(z)
    =
    y_{c^\star}^\star.
\]

This nearest-neighbor step can also be implemented as a secure computation black box. Existing work on privacy-preserving $k$-nearest-neighbor search and classification over encrypted or distributed data~\cite{wong2009secure,elmehdwi2014secure} shows that nearest-neighbor classification can be performed without revealing all private feature vectors or all distances.

In our setting, the required operation is simpler than full $k$-NN: we only need a secure nearest-neighbor search over the labeled representatives. The protocol should reveal the final predicted label, and possibly the identity of the nearest representative, but not the full set of distances.

\subsection{Pipeline Summary}

The full pipeline is:

\begin{enumerate}
    \item Use a black-box shared embedding method to map every private example to $\mathbb{R}^d$.
    \item Use a black-box private clustering method to assign every point to a cluster.
    \item For every client, cluster, and coordinate axis, create sorted local coordinate-value lists.
    \item Aggregate the sorted coordinate values so that the global sorted list $V_{c,\ell}$ is available for every cluster $c$ and coordinate $\ell$.
    \item Each client computes, for every local point, the exact sum of $L_1$ distances to all points in its cluster using the global sorted coordinate lists.
    \item For each cluster, use secure argmin to select the point with minimum total $L_1$ distance.
    \item Label the selected representatives.
    \item Classify future points using federated nearest neighbor over the labeled representatives.
\end{enumerate}

\end{document}
